\documentclass[10pt,a4paper]{amsart}
\usepackage[margin=19mm]{geometry}
\usepackage{amsmath,amssymb,mathtools}
\usepackage[scr=rsfs,cal=euler]{mathalfa}
\usepackage{xfrac}
\usepackage[unicode,hidelinks,bookmarks=false]{hyperref}
\newcommand{\ie}{i.e.\ }
\newcommand{\cf}{cf.\ }
\newcommand{\resp}{respectively\ }
\newcommand{\Subsection}{\subsection}
\providecommand{\nobreakdash}{\nobreak\hbox{-}}
\setlength{\emergencystretch}{2em}
\multlinegap0pt
\usepackage{bm}
\usepackage{bbm}
\usepackage{dsfont}
\usepackage{xspace}
\usepackage{cancel}
\usepackage{mathtools}
\usepackage{amsthm}
\makeatletter
\@ifundefined{theorem}{\newtheorem{theorem}{Theorem}}{}
\@ifundefined{lemma}{\newtheorem{lemma}[theorem]{Lemma}}{}
\makeatother
\newcommand{\cminus}{C^{-}}
\title[A Negative Answer to Stanley's Problem 4 on Differential Pos]{A Negative Answer to Stanley's Problem 4 on Differential Posets}
\author{Xinan Dai, Wenhao Deng, Yingdong Shi, Tailin Wu, Yuchen Yang}
\date{Source: arXiv:2607.24541v2 (CC BY 4.0)}
\begin{document}
\maketitle

\noindent\emph{Source and licence.} A Negative Answer to Stanley's Problem 4 on Differential Posets. Version arXiv:2607.24541v2 (CC BY 4.0), distributed under the Creative Commons Attribution 4.0 International licence, as recorded in the arXiv licence metadata for this exact version and shown on the abstract page https://arxiv.org/abs/2607.24541v2. Full rights notice: licenses/arxiv-2607.24541-CC-BY-4.0.txt.\par\smallskip

\noindent\textit{Editorial scope.} Counted unit: the Lemma whose source label is \texttt{lem:persistent} in the article (source0.tex lines 169--171, source label \texttt{lem:persistent}), delivered together with the article's own complete original proof of it (source0.tex lines 173--193, one \texttt{proof} environment of 21 lines). No mathematics is written, repaired or re-derived: statement and proof are the article's own text line for line. Required context retained uncounted: lem:crown (lines 148--150). Every definition, display and label of those lines is retained unchanged. Omitted from this package and not redistributed: 1--147, 151--168, 172--172, 194--346. Nothing inside a retained excerpt is omitted or altered: every retained body is delivered exactly as the article's lines are written, so no omission receipt of any kind accompanies this package. Citations: the retained excerpts carry no citation command at all, so no bibliography entry of the article is transcribed and no reference is redistributed. Reference adaptation: none was needed and none was made; every reference inside the delivered unit resolves inside it, so no unresolved reference or citation is produced. Printed slips kept literal: no printed word, symbol, hypothesis or number of the counted statement or of its proof was altered. Layout and class: the article's own class is \texttt{amsart} with packages including fontenc, amsmath, amssymb, mathtools, microtype, hyperref; the wrapper is the lane's pinned \texttt{amsart} editorial wrapper at 10pt on A4, which restates the theorem-like environments the retained text needs and adds the article's own macro definitions that the retained text needs (\texttt{\textbackslash cminus}), with extra wrapper packages (bm, bbm, dsfont, xspace, cancel, mathtools). The delivered numbering is the unit's own. Assets: no retained span contains a figure environment, a graphics-inclusion command, a PDF-page inclusion, a table or a TikZ picture; no image file is copied into this package, the package declares no asset dependency and no image file is redistributed with it. Licence: version arXiv:2607.24541v2 of this article is distributed under the Creative Commons Attribution 4.0 International licence, as recorded in the arXiv licence metadata for this exact version and shown on the abstract page https://arxiv.org/abs/2607.24541v2. A full rights notice is delivered at licenses/arxiv-2607.24541-CC-BY-4.0.txt. Characters: every retained excerpt is pure ASCII, so the delivered engine, XeLaTeX with its default Latin Modern text font, reports no missing character.
\begin{lemma}[Crown trade]\label{lem:crown}
Suppose a finite $1$-differential prefix through rank $n$ contains an element $y\in P_{n-1}$ with exactly two lower covers. Then the prefix has two extensions through rank $n+1$ whose new rank sizes differ by one.
\end{lemma}

\begin{lemma}[A persistent crownable element]\label{lem:persistent}
Every finite $1$-differential prefix with top rank at least four contains an element in the penultimate rank having exactly two lower covers and three upper covers. After either extension in Lemma~\ref{lem:crown}, the resulting prefix has the same property.
\end{lemma}

\begin{proof}
The ranks through $P_2$ are forced. Let $a_0=\widehat{0}$, let $a_1$ be the unique element of rank one, and denote the two elements of rank two by $a_2$ and $b_2$. We claim that every subsequent valid extension contains elements $a_i,b_i$ for which
\begin{equation}
\cminus(a_{i+1})=\{a_i\},
\qquad
\cminus(b_{i+1})=\{a_i,b_i\}
\qquad (i\geq2).
\label{eq:threads}
\end{equation}
Assume these elements have been chosen through rank $i$. The element $a_{i-1}$ has one lower cover, and therefore exactly two upper covers, namely $a_i$ and $b_i$. Since $a_i$ and $b_i$ share the lower cover $a_{i-1}$, condition (D2) gives a unique common upper cover $x$ in rank $i+1$. The element $a_i$ has one lower cover, so it has one further upper cover $z$.

We show that no additional element lies below either $x$ or $z$. If $w$ were such an element, then $a_i$ and $w$ would share an upper cover. By (D2), they would also share a lower cover. The only possibility is $a_{i-1}$, whose upper covers are $a_i$ and $b_i$. Thus $w=b_i$. This is already one of the lower covers of $x$, while for $z$ it would give a second common upper cover of $a_i$ and $b_i$. Both possibilities contradict the choice of $w$ as an additional lower cover. Hence
\[
\cminus(x)=\{a_i,b_i\},
\qquad
\cminus(z)=\{a_i\}.
\]
Taking $b_{i+1}=x$ and $a_{i+1}=z$ proves \eqref{eq:threads} inductively.

Now let the top rank be $P_n$ with $n\geq4$. The element $b_{n-1}$ has exactly two lower covers and, by (D3), three upper covers in $P_n$. It therefore supports the crown trade. Since the preceding argument applies to every valid one-rank extension, the required configuration is present again after either choice.
\end{proof}

\end{document}
